\chapter{Generador del Dataset v6 (\texttt{dataset\_generator.py})}
\label{anx:dataset-v6}

Este anexo contiene el código fuente completo del script \texttt{dataset\_generator.py}, responsable de sintetizar de manera determinística y parametrizada los pares de entrenamiento, validación y prueba (\textit{Out-Of-Distribution}) utilizados en el ajuste fino de CodeT5 (Dataset v6).

\begin{lstlisting}[caption={Script de generación programática del dataset v6 para CodeT5}, label={lst:dataset-generator-py}, numbers=left, breaklines=true, basicstyle=\ttfamily\small]
import json
import keyword
import random
from collections import Counter
from pathlib import Path

# Configuracion
SEED = 42
N_TRAIN = 200
N_VAL = 50
N_TEST_OOD = 50

OUTPUT_DIR = Path(__file__).resolve().parent / "dataset"

# Debe usarse identica en entrenamiento e inferencia
INSTRUCTION = "Generate Django model fields from the entity attributes."

# Tipos aceptados -> campo Django
DJANGO_FIELDS = {
    "string":  "models.CharField(max_length=255)",
    "int":     "models.IntegerField()",
    "float":   "models.FloatField()",
    "boolean": "models.BooleanField(default=False)",
    "date":    "models.DateField()",
    "time":    "models.TimeField()",
    "decimal": "models.DecimalField(max_digits=10, decimal_places=2)",
    "binary":  "models.BinaryField()",
}
TYPES = list(DJANGO_FIELDS)

# Pools de nombres (train vs held-out)
ENTITY_TRAIN = [
    "User", "Country", "City", "Company", "Department", "Student", "Teacher",
    "Book", "Author", "Category", "Order", "Product", "Invoice", "Customer",
    "Phone", "Role", "Address", "Payment", "Employee", "Project", "Task",
    "Course", "Subject", "Grade", "Supplier", "Warehouse", "Shipment",
    "Review", "Ticket", "Event",
    "Pais", "Ciudad", "Estudiante", "Docente", "Libro", "Autor", "Pedido",
    "Producto", "Factura", "Cliente", "Telefono", "Rol", "Usuario", "Empresa",
]
ENTITY_HELD = [
    "Hospital", "Paciente", "Doctor", "Cita", "Vehiculo", "Conductor",
    "Hotel", "Reserva", "Loan", "Library", "Sensor", "Device", "Matricula",
    "Sucursal", "Inventario", "Contrato", "Proveedor", "Receta", "Turno",
    "Mensaje",
]

ATTR_TRAIN = [
    "name", "email", "description", "title", "age", "stock", "quantity",
    "price", "salary", "is_active", "enabled", "created_at", "birth_date",
    "number", "status", "type", "code", "address", "phone", "city", "country",
    "carrera", "nombre", "edad", "precio", "fecha", "estado", "timestamp",
    "amount", "total", "score", "rating", "weight", "height", "start_time",
    "end_time", "photo", "file", "content", "role",
]
ATTR_HELD = [
    "apellido", "telefono", "direccion", "saldo", "descuento", "peso",
    "altura", "hora_inicio", "hora_fin", "avatar", "firma", "documento",
    "matricula", "semestre", "turno", "expires_at", "last_login", "nickname",
    "slug", "version",
]

SUFFIXES = ["id", "Id", "code", "ref", "value", "1", "2"]

# Validacion de nombres (la capa deterministica sanea el resto)
def valid_name(name):
    return (
        name.isidentifier()
        and name.isascii()
        and not keyword.iskeyword(name)
        and name != "pk"
        and "__" not in name
        and not name.startswith("_")
        and not name.endswith("_")
    )

# Generadores de nombres
def make_attr_name(rng, pool):
    base = rng.choice(pool)
    other = rng.choice(pool)
    r = rng.random()
    if r < 0.60:
        name = base
    elif r < 0.75:
        name = f"{base}_{rng.choice(SUFFIXES)}"
    elif r < 0.88:
        name = base + other.capitalize()
    else:
        name = f"{base}_{other}"
    return name


def make_entity_name(rng, pool):
    base = rng.choice(pool)
    other = rng.choice(pool)
    r = rng.random()
    if r < 0.65:
        name = base
    elif r < 0.85:
        name = f"{base}_{other.capitalize()}"
    else:
        name = base + other.lower()
    return name

# Generacion de un ejemplo
def make_entity(rng, entity_pool, attr_pool):
    while True:
        name = make_entity_name(rng, entity_pool)
        if valid_name(name):
            break

    amount = rng.choices(range(1, 9), weights=[12, 18, 20, 16, 12, 10, 6, 6])[0]

    attributes, used = [], set()
    while len(attributes) < amount:
        attr = make_attr_name(rng, attr_pool)
        if not valid_name(attr) or attr in used:
            continue
        used.add(attr)
        attributes.append({"name": attr, "type": rng.choice(TYPES)})

    return {"name": name, "attributes": attributes}


def build_sample(entity):
    model_input = json.dumps(entity, ensure_ascii=False, separators=(",", ":"))
    output = "\n".join(
        f"{a['name']} = {DJANGO_FIELDS[a['type']]}" for a in entity["attributes"]
    )
    return {"instruction": INSTRUCTION, "input": model_input, "output": output}


def generate_split(rng, amount, entity_pool, attr_pool, seen):
    samples = []
    attempts = 0
    while len(samples) < amount:
        attempts += 1
        if attempts > amount * 50:
            raise RuntimeError("No hay suficientes combinaciones unicas; amplia los pools.")
        sample = build_sample(make_entity(rng, entity_pool, attr_pool))
        if sample["input"] in seen:
            continue
        seen.add(sample["input"])
        samples.append(sample)
    return samples


def write_jsonl(path, samples):
    with open(path, "w", encoding="utf-8") as f:
        for s in samples:
            f.write(json.dumps(s, ensure_ascii=False) + "\n")

# Verificacion de consistencia
def verify(samples):
    for s in samples:
        entity = json.loads(s["input"])
        lines = s["output"].split("\n")
        assert len(lines) == len(entity["attributes"])
        for line, attr in zip(lines, entity["attributes"]):
            assert line.startswith(attr["name"] + " = models.")

# Main
def main():
    rng = random.Random(SEED)
    OUTPUT_DIR.mkdir(parents=True, exist_ok=True)

    seen = set()
    train = generate_split(rng, N_TRAIN, ENTITY_TRAIN, ATTR_TRAIN, seen)
    val = generate_split(rng, N_VAL, ENTITY_TRAIN, ATTR_TRAIN, seen)
    test_ood = generate_split(rng, N_TEST_OOD, ENTITY_HELD, ATTR_HELD, seen)

    for name, data in (("train", train), ("val", val), ("test_ood", test_ood)):
        verify(data)
        write_jsonl(OUTPUT_DIR / f"{name}.jsonl", data)

    types_seen = Counter(
        a["type"] for s in train for a in json.loads(s["input"])["attributes"]
    )
    print("Dataset v6 generado en:", OUTPUT_DIR)
    print(f"  train: {len(train)} | val: {len(val)} | test_ood: {len(test_ood)}")
    print("  distribucion de tipos (train):", dict(types_seen))
    print("  instruccion:", INSTRUCTION)


if __name__ == "__main__":
    main()
\end{lstlisting}

\chapter{Módulo de Restricciones Compuestas (\texttt{idiom\_constraints.py})}
\label{anx:idiom-constraints}

En este anexo se presenta la implementación de la clase personalizada \textit{CompositeForeignKeyConstraint}, la cual extiende \textit{BaseConstraint} para permitir la generación y gestión de restricciones de claves foráneas compuestas directamente en el esquema relacional de Django.

\begin{lstlisting}[caption={Implementación de CompositeForeignKeyConstraint para claves foráneas compuestas}, label={anx:code-idiom-constraints}]
from django.db.models.constraints import BaseConstraint


class CompositeForeignKeyConstraint(BaseConstraint):
    def __init__(
        self,
        *,
        name,
        fields,
        to_model,           # "app_label.NombreModelo"
        to_fields,
        on_delete="CASCADE",
        violation_error_code=None,
        violation_error_message=None,
    ):
        super().__init__(
            name=name,
            violation_error_code=violation_error_code,
            violation_error_message=violation_error_message,
        )
        self.fields = tuple(fields)
        self.to_model = to_model
        self.to_fields = tuple(to_fields)
        self.on_delete = on_delete

    def _clause(self, model, schema_editor):
        qn = schema_editor.quote_name
        target = model._meta.apps.get_model(self.to_model)
        cols = ", ".join(qn(model._meta.get_field(f).column) for f in self.fields)
        tcols = ", ".join(qn(target._meta.get_field(f).column) for f in self.to_fields)
        return (
            f"FOREIGN KEY ({cols}) "
            f"REFERENCES {qn(target._meta.db_table)} ({tcols}) "
            f"ON DELETE {self.on_delete} DEFERRABLE INITIALLY DEFERRED"
        )

    # Se usa cuando la restriccion va dentro del CREATE TABLE
    def constraint_sql(self, model, schema_editor):
        qn = schema_editor.quote_name
        return f"CONSTRAINT {qn(self.name)} {self._clause(model, schema_editor)}"

    # Se usa cuando se anade con ALTER TABLE (AddConstraint)
    def create_sql(self, model, schema_editor):
        qn = schema_editor.quote_name
        return (
            f"ALTER TABLE {qn(model._meta.db_table)} "
            f"ADD CONSTRAINT {qn(self.name)} {self._clause(model, schema_editor)}"
        )

    def remove_sql(self, model, schema_editor):
        qn = schema_editor.quote_name
        return f"ALTER TABLE {qn(model._meta.db_table)} DROP CONSTRAINT {qn(self.name)}"

    def validate(self, model, instance, exclude=None, using="default"):
        pass  # la integridad la garantiza la base de datos

    def __eq__(self, other):
        if isinstance(other, CompositeForeignKeyConstraint):
            return (
                self.name == other.name
                and self.fields == other.fields
                and self.to_model == other.to_model
                and self.to_fields == other.to_fields
                and self.on_delete == other.on_delete
            )
        return super().__eq__(other)

    def __hash__(self):
        return hash((self.name, self.fields, self.to_model, self.to_fields, self.on_delete))

    def __repr__(self):
        return f"<CompositeForeignKeyConstraint name={self.name!r} {self.fields} -> {self.to_model}{self.to_fields}>"

    def deconstruct(self):
        path, args, kwargs = super().deconstruct()
        kwargs.update(
            fields=self.fields,
            to_model=self.to_model,
            to_fields=self.to_fields,
            on_delete=self.on_delete,
        )
        return path, args, kwargs
\end{lstlisting}

\chapter{Módulo de Secuencia Concurrente (\texttt{idiom\_sequences.py})}
\label{anx:idiom-sequences}

En este anexo se presenta el código fuente del \textit{mixin} desarrollado para gestionar la asignación segura de identificadores consecutivos en modelos con claves primarias complejas bajo condiciones de concurrencia.

\begin{lstlisting}[caption={Implementación del SequentialIdMixin para el manejo de bloqueos y secuencias atómicas}, label={anx:code-idiom-sequences}]
from django.db import connections, router, transaction
from django.db.models import Max
class SequentialIdMixin:
    sequential_field: str = ""
    def save(self, *args, **kwargs):
        if args:
            raise TypeError("save() solo admite argumentos con nombre")
        name = self.sequential_field
        if not name:
            raise ValueError(f"{type(self).__name__} debe definir 'sequential_field'")

        if getattr(self, name) is not None:
            return super().save(**kwargs)

        using = kwargs.get("using") or router.db_for_write(type(self), instance=self)
        with transaction.atomic(using=using):
            connection = connections[using]
            table = connection.ops.quote_name(self._meta.db_table)
            with connection.cursor() as cursor:
                cursor.execute(f"LOCK TABLE {table} IN EXCLUSIVE MODE")
            current = type(self)._default_manager.using(using).aggregate(m=Max(name))["m"] or 0
            setattr(self, name, current + 1)
            kwargs["force_insert"] = True
            return super().save(**kwargs)
\end{lstlisting}